\section{Round-nine reconstruction: constrained curvature, covariance-first tangents, and kinetic observability}
\label{sec:r9-b3}

The raw second variation of the perspective entropy is not asserted to be
positive.  Positivity is obtained only on a regular, sufficiently short
constrained phase by combining the dynamic Hessian with the multiplier
curvature dictated by the linearized balance.  Independently, finite-volume
cumulants first construct the Gaussian covariance.  The variational quadratic
form is then identified with that covariance by a second-order epi-convergence
theorem, avoiding the circular definition of one through the other.

\subsection{Weighted tangent and cotangent spaces}

Fix $m>8$.  Let $\mathsf X_T$ be the completion of smooth density tangents for
\[
 \|u\|_{\mathsf X_T}^2
 =\sup_{t\le T}\int\frac{|u_t|^2}{f_t}(1+|v|^m)
 +\int_0^T\|u_t\|_{H^{-1}_{x,v}(w_m)}^2dt,
\]
and let $\mathsf Y_T$ be the weighted $L^2((qA_f)^{-1})$ space of contact
current tangents.  The closed linearized balance operator is
\[
 \mathcal B(u,J)
 =\partial_tu+v\cdot\nabla_xu-\Delta^*J.
\]
Together with its endpoint and conservation coordinates.  Its kernel with the
microcanonical initial tangent condition is denoted $\mathsf T_{f,q}$.
The cotangent pair $(p,\psi)$ is taken in the corresponding weighted transport
Sobolev/Orlicz space and acts through $\Delta p+\psi$.

For a regular biased path $(f,qA_f)$ the exact unconstrained second variation
is
\[
\begin{aligned}
 \mathfrak h_{f,q}(u,J)
 ={}&D^2I_0(f_0)[u_0,u_0]
 +\int_0^T\!\int
 \frac{|J-qDA_f[u]|^2}{q}\,dA_fdt\\
 &+\int_0^T\!\int(1-q)D^2A_f[u,u]dt.
\end{aligned}
\]
The last term can have either sign.

\subsection{The KKT-completed tangent curvature}

Let $p^*$ be the regular Lagrange multiplier for the weak balance at the
biased phase.  The constrained Lagrangian is
\[
 \mathscr L(f,\Gamma,p)
 =I_0(f_0)+\int\ell(d\Gamma/dA_f)dA_f
 +\langle p,\mathcal B(f,\Gamma)\rangle.
\]
Its second variation in $(f,\Gamma)$ contains the additional multiplier
curvature
\[
 \mathfrak c_{p^*}(u)
 =\int_0^T\!\int qA_f\,
  D^2_f(\Delta p^*)[u,u]dt
 +\text{endpoint preparation curvature}.
\]
For the binary Boltzmann reference the first term is the polarization of the
bilinear collision intensity evaluated against the biased collision
increment.

\begin{lemma}[Linearized kinetic energy estimate]
\label{lem:r9-b3-energy}
On every compact regular phase chart there are $C_0,C_1$ such that every
$(u,J)\in\mathsf T_{f,q}$ satisfies
\[
 \|u\|_{\mathsf X_T}^2
 \le C_0\|u_0\|_{L^2(f_0^{-1}w_m)}^2
 +C_1\int_0^T\!\int
 \frac{|J-qDA_f[u]|^2}{q}\,dA_fdt.
\]
The constants are uniform for $T\le T_*$.  Moreover
\[
 |\mathfrak h_{f,q}^{\rm neg}(u)+\mathfrak c_{p^*}^{\rm neg}(u)|
 \le CT\|q-1\|_{L^\infty_w}\|u\|_{\mathsf X_T}^2.
\]
\end{lemma}

\begin{proof}
Test the linearized balance with the solution of the backward weighted
transport--collision equation having datum $u/f$.  Free transport is
skew-adjoint after the endpoint terms, while the symmetric part of the
linearized collision operator controls the orthogonal complement of the
collision invariants.  The finite-dimensional invariant modes are fixed by
the initial microcanonical tangent condition.  Cauchy--Schwarz gives the
contact residual term and Gronwall gives the first estimate.  The bilinearity
of $A_f$, boundedness of the regular multiplier, and the velocity moments give
the second estimate.  All boundary pairings are justified by the B2 closed
Green graph.
\end{proof}

\begin{theorem}[Positive constrained curvature]
\label{thm:r9-b3-curvature}
Choose the regular horizon so that
\[
 CT_*\|q-1\|_{L^\infty_w}<\frac12.
\]
Then the KKT-completed form
\[
 \mathfrak q_{f,q}(u,J)
 =\mathfrak h_{f,q}(u,J)+\mathfrak c_{p^*}(u)
\]
is continuous and coercive on $\mathsf T_{f,q}$:
\[
 \mathfrak q_{f,q}(u,J)
 \ge c\left(
  \|u_0\|_{L^2(f_0^{-1}w_m)}^2+
  \|u\|_{\mathsf X_T}^2+
  \int\frac{|J-qDA_f[u]|^2}{q}dA_fdt
 \right).
\]
No positivity is claimed for the raw perspective Hessian off the constrained
tangent or outside the regular short-time phase chart.
\end{theorem}

\begin{proof}
The positive initial entropy and square contact residual control the right
side of Lemma~\ref{lem:r9-b3-energy}.  Absorb the combined negative curvature
using the displayed small-horizon condition.  The remaining KKT terms are
bounded symmetric perturbations, proving continuity and coercivity.
\end{proof}

\subsection{Closed range and the complete gauge}

Let $\mathcal B^*$ be the backward kinetic adjoint.  Constants and collision
invariants form a finite-dimensional space $\mathcal N$; fix their endpoint
normalization.

\begin{theorem}[Backward observability and closed balance range]
\label{thm:r9-b3-range}
For every normalized cotangent $(p,\psi)$,
\[
 \|(p,\psi)\|_{\mathsf P_T/\mathcal N}
 \le C\left(
  \|\mathcal B^*(p,\psi)\|_{\mathsf T_{f,q}^*}
  +\|\Delta p+\psi\|_{L^2(qA_f)}
 \right).
\]
Consequently $\mathcal B$ has closed range and
\[
 \ker\mathfrak D
 =\{(r,-\Delta r):r\in\mathsf P_T\},
 \qquad
 \mathfrak D(p,\psi)=\Delta p+\psi.
\]
Thus the identifiable cotangent is the Hausdorff quotient by the full balance
gauge, not merely by collision invariants.
\end{theorem}

\begin{proof}
Solve the backward equation blockwise by Duhamel around the free-transport
characteristics.  The symmetric collision part has a spectral gap off
$\mathcal N$; the macroscopic modes satisfy a finite hyperbolic system whose
endpoint data are controlled by the balance functional.  Velocity truncation
and the B2 trace estimate pass the result to the weighted spaces.  This gives
the observability inequality.  The Hilbert closed-range theorem applies to
$\mathcal B$.  If $\mathfrak D(p,\psi)=0$, the dual action on every balanced
tangent vanishes; closed range and the adjoint characterization give
$(p,\psi)=(r,-\Delta r)$ after the invariant normalization.  The converse is
the exact weak balance identity.
\end{proof}

\subsection{Finite-volume covariance first}

Let $Z_\varepsilon$ be the exact-centred joint density/contact fluctuation
field under a regular finite-volume source:
\[
 Z_\varepsilon
 =\sqrt{\mu_\varepsilon}
 \left((\pi^\varepsilon,\Gamma^\varepsilon)
 -D Q_\varepsilon(\Theta)\right).
\]
For a finite test family $F=(F_1,\ldots,F_d)$ set
$\Sigma_\varepsilon(F)=D^2Q_\varepsilon(\Theta)[F,F]$.

\begin{theorem}[Finite-dimensional covariance and Gaussian limit]
\label{thm:r9-b3-finiteclt}
For every finite weighted test family,
\[
 \Sigma_\varepsilon(F)\longrightarrow\Sigma(F),
\]
and the characteristic function of $F(Z_\varepsilon)$ converges to that of
$N(0,\Sigma(F))$.  The family $(\Sigma(F))$ is projectively compatible and
positive semidefinite.  Its null space consists precisely of the balance gauge
and the prepared conservation directions.
\end{theorem}

\begin{proof}
Normal convergence of the B2 complex marked pressure gives convergence of all
finite derivatives.  The $k$th cumulant of the scaled field is
$O(\mu_\varepsilon^{1-k/2})$ for $k\ge3$, while the second converges to the
Hessian.  This proves the Gaussian characteristic function.  Positivity is
inherited from finite covariance matrices.  A zero-variance class is constant
under every local activity/contact perturbation.  Varying one free-flight cell
and one actual binary contact shows that it annihilates every balanced tangent;
Theorem~\ref{thm:r9-b3-range} identifies it with the full gauge plus prepared
constraints.
\end{proof}

The compatible finite-dimensional covariances define a Gaussian cylindrical
measure on the dual of the nuclear test core.  They do not imply a bounded
inverse on the whole Hilbert completion.

\subsection{Second-order epi convergence}

Let $I_\varepsilon$ be the finite-volume convex rate on a finite projection,
let $x_\varepsilon=DQ_\varepsilon(\Theta)$, and define
\[
 J_\varepsilon(z)=
 \mu_\varepsilon\left[
 I_\varepsilon(x_\varepsilon+z/\sqrt{\mu_\varepsilon})
 -I_\varepsilon(x_\varepsilon)
 -\Theta\cdot z/\sqrt{\mu_\varepsilon}
 \right].
\]

\begin{theorem}[Mosco identification of the tangent action]
\label{thm:r9-b3-mosco}
On every finite projection, $J_\varepsilon$ Mosco-converges to
$\frac12\Sigma^{-1}$ on the covariance range.  Passing through the increasing
balanced test core yields the closed quadratic form associated with
$\mathfrak q_{f,q}$ from Theorem~\ref{thm:r9-b3-curvature}.  Hence the
Lax--Milgram inverse of the KKT-completed form and the covariance constructed
in Theorem~\ref{thm:r9-b3-finiteclt} agree on their common form domain.
\end{theorem}

\begin{proof}
On a finite projection, analytic Taylor expansion of $Q_\varepsilon$ and
Legendre duality give the quadratic liminf.  For a vector in the covariance
range, tilt by $\Theta+h/\sqrt{\mu_\varepsilon}$ to obtain a recovery sequence
and the matching limsup.  Projectively increasing the finite test spaces gives
a monotone sequence of closed forms.  The balance constraint and B1 initial
Schur complement are continuous on the form domain.  The explicit second
variation of the constrained Lagrangian is exactly
$\mathfrak q_{f,q}$, so uniqueness of the Mosco limit identifies the two
forms.  Coercivity on the regular chart then gives the Lax--Milgram
representation.
\end{proof}

\subsection{Process tightness from localized cumulants}

For a time interval $I$ and a test $F$, let $Z_\varepsilon(I,F)$ be the
increment of the exact-centred field.  The B2 source expansion localized to
deterministic intervals gives connected cumulant densities with one anchor
time and tree-decaying relative times.

\begin{lemma}[Fourth-moment increment estimate]
\label{lem:r9-b3-fourth}
For every finite test family there are $C,\eta>0$ such that
\[
 \mathbb E|Z_\varepsilon([s,t],F)|^4
 \le C|t-s|^{2}+C\mu_\varepsilon^{-1}|t-s|,
\]
and
\[
 \mathbb E|Z_\varepsilon([s,t],F)|^2
 \le C|t-s|.
\]
The bounds include intervals sharing an endpoint; no product
$|I||J|$ estimate is assumed without a cumulant-density integration.
\end{lemma}

\begin{proof}
Differentiate the B2 pressure with four independently supported time sources.
A connected tree has one anchor time and each additional relative event time
is integrated against the renewal kernel.  Its $L^1$ norm is summable on the
short horizon.  Integrating the four-point cumulant density over
$[s,t]^4$ gives $C|t-s|^2$; diagonal coincidences contribute the finite-size
term $C\mu^{-1}|t-s|$.  The two-point statement is identical.  Boundary atoms
are already part of the trace-source cumulant and need no stopping-time
substitution.
\end{proof}

\begin{theorem}[Prepared joint Gaussian process]
\label{thm:r9-b3-main}
The density/contact fluctuation fields converge in the weighted distribution
Skorokhod space to the unique centred Gaussian process with covariance
$\Sigma$ and linearized biased balance.  The initial covariance is the B1
prepared Schur complement.  Its Cameron--Martin form is the closed
KKT-completed form of Theorem~\ref{thm:r9-b3-mosco}.
\end{theorem}

\begin{proof}
Finite-dimensional convergence is
Theorem~\ref{thm:r9-b3-finiteclt}.  Lemma~\ref{lem:r9-b3-fourth} gives
Kolmogorov/Aldous bounds for each member of a countable nuclear test core;
Mitoma's criterion gives process tightness.  Passing the exact microscopic
Green balance to the limit identifies the linearized martingale problem.
The covariance and initial law determine its unique Gaussian solution.
Theorem~\ref{thm:r9-b3-mosco} identifies the Cameron--Martin form.
\end{proof}
